The first data-mining pass evaluates whether the available data can support the recoverable-resilience research program before the optimization model is introduced. The analysis covers data availability, rainfall shock characterization, speed-deficit and recovery proxies, spatial concentration of deficit burden, demand-network structure, and alignment between preliminary resilience metrics and network exposure.

Across the usable speed datasets, the strongest observed functional-loss signals appear in Chicago and New York. Chicago has a high p90 speed deficit and a large share of observations with at least 20 percent speed deficit, while New York has a weaker but still substantial deficit signal over a much larger observed TMC set. Rainfall-speed alignment is positive but moderate: San Antonio, Austin, New York, and Chicago show the clearest lagged association between rainfall and speed deficit, with event-level recovery proxies typically on the order of a few hours in the current monthly data.

The spatial concentration results are important for the paper's decision-centered framing. In several cities, the top 10 percent of TMCs account for roughly 30--40 percent of the total speed-deficit burden. This does not prove recoverability, but it shows that observed loss is not uniformly diffuse. A targeted intervention model may therefore have meaningful spatial leverage to test.

Demand and network data provide a baseline empirical foundation for functional dependence. All 11 cities have OD demand and link-performance outputs, with strong cross-city differences in OD scale, destination concentration, and congested-volume share. These data can support a baseline dependence matrix \(Q_t\), although destination service importance \(S_j\) still requires additional POI, employment, health-care, retail, population, or other functional-exposure data.

Finally, for the three cities with existing resilience-index files, the observed loss metrics can be matched to link-performance exposure. Chicago shows the clearest positive alignment between loss magnitude and link volume; New York shows weaker but positive alignment; Houston is weak or mixed. This suggests that the data contain functional-exposure signal, but descriptive alignment alone cannot answer how much loss would be recoverable under alternative intervention allocations.
